% ============================================================
% 1. INTRODUCTION
% ============================================================
\section{Introduction}
\label{sec:introduction}

Medical imaging plays a central role in the diagnosis and monitoring of neurological disorders, with Magnetic Resonance Imaging (MRI) serving as the primary modality for detecting structural brain abnormalities such as tumors, lesions, and atrophy~\cite{baur2021autoencoders}. Despite its diagnostic value, the clinical deployment of automated analysis systems faces persistent challenges: annotated medical datasets are scarce due to the cost and expertise required for labeling, class distributions are heavily imbalanced (pathological cases are far less common than healthy ones), and privacy regulations constrain data sharing across institutions~\cite{yi2019generative}.

Traditional supervised approaches to brain tumor detection require large collections of pixel-level annotations, which are expensive and time-consuming to acquire. Furthermore, discriminative classifiers trained on limited datasets often fail to generalize across imaging protocols, scanner manufacturers, and patient demographics. These limitations have motivated growing interest in \textit{unsupervised} anomaly detection methods, which learn a model of normality exclusively from healthy data and identify pathology as deviations from this learned distribution~\cite{baur2021autoencoders, schlegl2019f}.

Generative models have emerged as a promising paradigm for both anomaly detection and data augmentation in medical imaging. Generative Adversarial Networks (GANs)~\cite{goodfellow2014generative} were among the first deep generative architectures applied to medical image synthesis~\cite{yi2019generative}, but they suffer from well-documented training instabilities, mode collapse, and limited diversity. More recently, Denoising Diffusion Probabilistic Models (DDPMs)~\cite{ho2020denoising} have demonstrated superior sample quality and training stability compared to GANs, achieving state-of-the-art results in natural image synthesis~\cite{dhariwal2021diffusion}. Latent Diffusion Models (LDMs)~\cite{rombach2022high} extend this paradigm by operating the diffusion process in a compressed latent space, dramatically reducing computational requirements while preserving generation quality.

In the context of medical anomaly detection, diffusion-based methods reconstruct pathological inputs as their healthy counterparts, and the resulting reconstruction error serves as an anomaly signal~\cite{wyatt2022anoddpm, wolleb2022diffusion, sanchez2022healthy}. This approach is particularly attractive because it requires \textit{no anomaly labels during training}---the model learns only from healthy examples and can detect arbitrary pathologies at inference time.

In this work, we present a complete latent diffusion pipeline for brain MRI analysis that serves a dual purpose: (1)~anomaly detection via reconstruction-error-based tumor localization, and (2)~high-quality synthetic brain MRI generation for potential data augmentation. Our contributions are as follows:

\begin{enumerate}[leftmargin=*]
  \item We design and implement a two-stage latent diffusion architecture combining a MONAI-based Variational Autoencoder (VAE) with a class-conditional U-Net DDPM, featuring Classifier-Free Guidance (CFG) and DDIM sampling for efficient inference.
  \item We conduct a systematic hyperparameter ablation study across 120 configurations---spanning 10~model checkpoints, 4~noise injection levels, and 3~guidance scales---to characterize the sensitivity of anomaly detection performance to each factor.
  \item We provide quantitative evaluation of both anomaly detection (image-level and pixel-level AUROC) and synthetic image quality (FID, SSIM, LPIPS, diversity).
  \item We incorporate explainability mechanisms through self-attention map visualization and denoising trajectory analysis, and discuss ethical considerations for clinical deployment.
\end{enumerate}

The remainder of this paper is organized as follows. Section~\ref{sec:related_work} reviews related work on diffusion models, medical anomaly detection, and synthetic image generation. Section~\ref{sec:methodology} details our two-stage architecture and inference pipelines. Section~\ref{sec:experiments} describes the experimental setup, dataset, and evaluation metrics. Section~\ref{sec:results} presents quantitative and qualitative results. Sections~\ref{sec:discussion} and~\ref{sec:ethics} provide analysis, limitations, and ethical considerations. Section~\ref{sec:conclusion} concludes the paper.
